\documentclass{amsart}
% For dealing with references we use the comment environment
\usepackage{verbatim}
\newenvironment{reference}{\comment}{\endcomment}
%\newenvironment{reference}{}{}
\newenvironment{slogan}{\comment}{\endcomment}
\newenvironment{history}{\comment}{\endcomment}

% For commutative diagrams we use Xy-pic
\usepackage[all]{xy}

% We use 2cell for 2-commutative diagrams.
\xyoption{2cell}
\UseAllTwocells

% We use multicol for the list of chapters between chapters
\usepackage{multicol}

% This is generally recommended for better output
\usepackage{lmodern}
\usepackage[T1]{fontenc}

% For cross-file-references

% Package for hypertext links:
\usepackage{hyperref}

% For any local file, say "hello.tex" you want to link to please

% Theorem environments.
%
\theoremstyle{plain}
\newtheorem{theorem}[subsection]{Theorem}
\newtheorem{proposition}[subsection]{Proposition}
\newtheorem{lemma}[subsection]{Lemma}

\theoremstyle{definition}
\newtheorem{definition}[subsection]{Definition}
\newtheorem{example}[subsection]{Example}
\newtheorem{exercise}[subsection]{Exercise}
\newtheorem{situation}[subsection]{Situation}

\theoremstyle{remark}
\newtheorem{remark}[subsection]{Remark}
\newtheorem{remarks}[subsection]{Remarks}

\numberwithin{equation}{subsection}

% Macros
%
\def\lim{\mathop{\mathrm{lim}}\nolimits}
\def\colim{\mathop{\mathrm{colim}}\nolimits}
\def\Spec{\mathop{\mathrm{Spec}}}
\def\Hom{\mathop{\mathrm{Hom}}\nolimits}
\def\Ext{\mathop{\mathrm{Ext}}\nolimits}
\def\SheafHom{\mathop{\mathcal{H}\!\mathit{om}}\nolimits}
\def\SheafExt{\mathop{\mathcal{E}\!\mathit{xt}}\nolimits}
\def\Sch{\mathit{Sch}}
\def\Mor{\mathop{\mathrm{Mor}}\nolimits}
\def\Ob{\mathop{\mathrm{Ob}}\nolimits}
\def\Sh{\mathop{\mathit{Sh}}\nolimits}
\def\NL{\mathop{N\!L}\nolimits}
\def\CH{\mathop{\mathrm{CH}}\nolimits}
\def\proetale{{pro\text{-}\acute{e}tale}}
\def\etale{{\acute{e}tale}}
\def\QCoh{\mathit{QCoh}}
\def\Ker{\mathop{\mathrm{Ker}}}
\def\Im{\mathop{\mathrm{Im}}}
\def\Coker{\mathop{\mathrm{Coker}}}
\def\Coim{\mathop{\mathrm{Coim}}}

% Boxtimes
%
\DeclareMathSymbol{\boxtimes}{\mathbin}{AMSa}{"02}

%
% Macros for moduli stacks/spaces
%
\def\QCohstack{\mathcal{QC}\!\mathit{oh}}
\def\Cohstack{\mathcal{C}\!\mathit{oh}}
\def\Spacesstack{\mathcal{S}\!\mathit{paces}}
\def\Quotfunctor{\mathrm{Quot}}
\def\Hilbfunctor{\mathrm{Hilb}}
\def\Curvesstack{\mathcal{C}\!\mathit{urves}}
\def\Polarizedstack{\mathcal{P}\!\mathit{olarized}}
\def\Complexesstack{\mathcal{C}\!\mathit{omplexes}}
% \Pic is the operator that assigns to X its picard group, usage \Pic(X)
% \Picardstack_{X/B} denotes the Picard stack of X over B
% \Picardfunctor_{X/B} denotes the Picard functor of X over B
\def\Pic{\mathop{\mathrm{Pic}}\nolimits}
\def\Picardstack{\mathcal{P}\!\mathit{ic}}
\def\Picardfunctor{\mathrm{Pic}}
\def\Deformationcategory{\mathcal{D}\!\mathit{ef}}

\setlength{\emergencystretch}{3em}
\begin{document}
\title[Stacks Project, Tag 04R2]{634 --- Separate etale locality yields source-and-target locality}
\maketitle
\noindent Copyright (C) 2005--2025 Johan de Jong. Stacks Project authors. Source: \href{https://stacks.math.columbia.edu/tag/04R2}{Tag 04R2}.

\noindent Editorial difficulty note (not author text). Difficulty H2: The proof reduces through three affine localization steps and then inducts on the universally bounded degrees of etale fibres. It removes the diagonal from a fibre product to lower this degree, a geometric mechanism substantially beyond qualification exercises..

\noindent Editorial provenance note (not author text). History: excerpt from revision \texttt{a0444\allowbreak 6e57e\allowbreak c1fbc\allowbreak 252a8\allowbreak 71afc\allowbreak ec775\allowbreak 2fb28\allowbreak 07b14}, descent.tex, lines 7314--7490. Reference presentation was adapted; original-author context and core support blocks were retained or added. The mathematical wording is unchanged. Licensed under GNU FDL 1.2 or later; no invariant sections or cover texts. See ../licenses/COPYING. Context blocks reproduce author definitions and supporting results; they are not additional counted problems.

\begin{definition}
\label{definition-local-source-target}
Let $\mathcal{P}$ be a property of morphisms of schemes.
We say $\mathcal{P}$ is {\it \'etale local on source-and-target} if
\begin{enumerate}
\item (stable under precomposing with \'etale maps)
if $f : X \to Y$ is \'etale and $g : Y \to Z$ has $\mathcal{P}$,
then $g \circ f$ has $\mathcal{P}$,
\item (stable under \'etale base change)
if $f : X \to Y$ has $\mathcal{P}$ and $Y' \to Y$ is \'etale, then
the base change $f' : Y' \times_Y X \to Y'$ has $\mathcal{P}$, and
\item (locality) given a morphism $f : X \to Y$ the following are equivalent
\begin{enumerate}
\item $f$ has $\mathcal{P}$,
\item for every $x \in X$ there exists a commutative diagram
$$
\xymatrix{
U \ar[d]_a \ar[r]_h & V \ar[d]^b \\
X \ar[r]^f & Y
}
$$
with \'etale vertical arrows and $u \in U$ with $a(u) = x$ such that
$h$ has $\mathcal{P}$.
\end{enumerate}
\end{enumerate}
\end{definition}

\begin{lemma}
\label{lemma-universally-injective-etale-open-immersion}
A universally injective \'etale morphism is an open immersion.
\end{lemma}

\begin{proof}[First proof]
Let $f : X \to Y$ be an \'etale morphism which is universally injective.
Then $f$ is open
(Morphisms, Lemma \href{https://stacks.math.columbia.edu/tag/03WT}{lemma etale open (Tag 03WT)})
hence we can replace $Y$ by $f(X)$ and we may assume that $f$ is surjective.
Then $f$ is bijective and open hence a homeomorphism. Hence $f$ is
quasi-compact. Thus by
Lemma \href{https://stacks.math.columbia.edu/tag/06NC}{lemma flat surjective quasi compact monomorphism isomorphism (Tag 06NC)}
it suffices to show that $f$ is a monomorphism. As $X \to Y$ is \'etale
the morphism $\Delta_{X/Y} : X \to X \times_Y X$ is an open immersion by
Morphisms, Lemma \href{https://stacks.math.columbia.edu/tag/02GE}{lemma diagonal unramified morphism (Tag 02GE)}
(and
Morphisms, Lemma \href{https://stacks.math.columbia.edu/tag/02GV}{lemma flat unramified etale (Tag 02GV)}).
As $f$ is universally injective $\Delta_{X/Y}$ is also surjective, see
Morphisms, Lemma \href{https://stacks.math.columbia.edu/tag/01S4}{lemma universally injective (Tag 01S4)}.
Hence $\Delta_{X/Y}$ is an isomorphism, i.e., $X \to Y$ is a monomorphism.
\end{proof}

\begin{proof}[Second proof]
Let $f : X \to Y$ be an \'etale morphism which is universally injective.
Then $f$ is open (Morphisms, Lemma \href{https://stacks.math.columbia.edu/tag/03WT}{lemma etale open (Tag 03WT)})
hence we can replace $Y$ by $f(X)$ and we may assume that $f$ is surjective.
Since the hypotheses remain satisfied after any base change, we conclude
that $f$ is a universal homeomorphism. Therefore $f$ is integral, see
Morphisms, Lemma \href{https://stacks.math.columbia.edu/tag/04DF}{lemma universal homeomorphism (Tag 04DF)}.
It follows that $f$ is finite by
Morphisms, Lemma \href{https://stacks.math.columbia.edu/tag/01WJ}{lemma finite integral (Tag 01WJ)}.
It follows that $f$ is finite locally free by
Morphisms, Lemma \href{https://stacks.math.columbia.edu/tag/02KB}{lemma finite flat (Tag 02KB)}.
To finish the proof, it suffices that $f$ is finite locally
free of degree $1$ (a finite locally free morphism of degree $1$
is an isomorphism).
There is decomposition of $Y$ into open and closed subschemes
$V_d$ such that $f^{-1}(V_d) \to V_d$ is finite locally free of
degree $d$, see Morphisms, Lemma \href{https://stacks.math.columbia.edu/tag/04MH}{lemma finite locally free (Tag 04MH)}.
If $V_d$ is not empty, we can pick a morphism $\Spec(k) \to V_d \subset Y$
where $k$ is an algebraically closed field (just take the algebraic
closure of the residue field of some point of $V_d$).
Then $\Spec(k) \times_Y X \to \Spec(k)$ is a disjoint union of
copies of $\Spec(k)$, by
Morphisms, Lemma \href{https://stacks.math.columbia.edu/tag/02GL}{lemma etale over field (Tag 02GL)}
and the fact that $k$ is algebraically closed.
However, since $f$ is universally injective, there can only be
one copy and hence $d = 1$ as desired.
\end{proof}


\begin{lemma}
\label{lemma-etale-local-source-target}
Let $\mathcal{P}$ be a property of morphisms of schemes.
Assume
\begin{enumerate}
\item $\mathcal{P}$ is \'etale local on the source,
\item $\mathcal{P}$ is \'etale local on the target, and
\item $\mathcal{P}$ is stable under postcomposing with open immersions:
if $f : X \to Y$ has $\mathcal{P}$ and $Y \subset Z$ is an open
subscheme then $X \to Z$ has $\mathcal{P}$.
\end{enumerate}
Then $\mathcal{P}$ is \'etale local on the source-and-target.
\end{lemma}

\begin{proof}
Let $\mathcal{P}$ be a property of morphisms of schemes which
satisfies conditions (1), (2) and (3) of the lemma. By
Lemma \href{https://stacks.math.columbia.edu/tag/04QV}{lemma precompose property local source (Tag 04QV)}
we see that $\mathcal{P}$ is stable under precomposing with
\'etale morphisms. By
Lemma \href{https://stacks.math.columbia.edu/tag/04QU}{lemma pullback property local target (Tag 04QU)}
we see that $\mathcal{P}$ is stable under \'etale base change.
Hence it suffices to prove part (3) of
Definition \ref{definition-local-source-target}
holds.

\medskip\noindent
More precisely, suppose that $f : X \to Y$ is a morphism
of schemes which satisfies
Definition \ref{definition-local-source-target} part (3)(b).
In other words, for every $x \in X$ there exists an \'etale
morphism $a_x : U_x \to X$, a point $u_x \in U_x$ mapping to $x$,
an \'etale morphism $b_x : V_x \to Y$, and a morphism $h_x : U_x \to V_x$
such that $f \circ a_x = b_x \circ h_x$ and $h_x$ has $\mathcal{P}$.
The proof of the lemma is complete once we show that $f$ has $\mathcal{P}$.
Set $U = \coprod U_x$, $a = \coprod a_x$, $V = \coprod V_x$,
$b = \coprod b_x$, and $h = \coprod h_x$. We obtain a
commutative diagram
$$
\xymatrix{
U \ar[d]_a \ar[r]_h & V \ar[d]^b \\
X \ar[r]^f & Y
}
$$
with $a$, $b$ \'etale, $a$ surjective. Note that $h$ has $\mathcal{P}$
as each $h_x$ does and $\mathcal{P}$ is \'etale local on the target.
Because $a$ is surjective and $\mathcal{P}$ is \'etale local on the source,
it suffices to prove that $b \circ h$ has $\mathcal{P}$.
This reduces the lemma to proving that $\mathcal{P}$ is stable under
postcomposing with an \'etale morphism.

\medskip\noindent
During the rest of the proof we let $f : X \to Y$ be a
morphism with property $\mathcal{P}$ and $g : Y \to Z$ is an \'etale
morphism. Consider the following statements:
\begin{enumerate}
\item[(-)] With no additional assumptions $g \circ f$
has property $\mathcal{P}$.
\item[(A)] Whenever $Z$ is affine
$g \circ f$ has property $\mathcal{P}$.
\item[(AA)] Whenever $X$ and $Z$ are affine
$g \circ f$ has property $\mathcal{P}$.
\item[(AAA)] Whenever $X$, $Y$, and $Z$ are affine
$g \circ f$ has property $\mathcal{P}$.
\end{enumerate}
Once we have proved (-) the proof of the lemma will be complete.

\medskip\noindent
Claim 1: (AAA) $\Rightarrow$ (AA).
Namely, let $f : X \to Y$, $g : Y \to Z$ be as above with $X$, $Z$ affine.
As $X$ is affine hence quasi-compact we can find finitely many
affine open $Y_i \subset Y$, $i = 1, \ldots, n$ such that
$X = \bigcup_{i = 1, \ldots, n} f^{-1}(Y_i)$. Set $X_i = f^{-1}(Y_i)$. By
Lemma \href{https://stacks.math.columbia.edu/tag/04QU}{lemma pullback property local target (Tag 04QU)}
each of the morphisms $X_i \to Y_i$ has $\mathcal{P}$.
Hence $\coprod_{i = 1, \ldots, n} X_i \to \coprod_{i = 1, \ldots, n} Y_i$
has $\mathcal{P}$ as $\mathcal{P}$ is \'etale local on the target.
By (AAA) applied to
$\coprod_{i = 1, \ldots, n} X_i \to \coprod_{i = 1, \ldots, n} Y_i$
and the \'etale morphism $\coprod_{i = 1, \ldots, n} Y_i \to Z$
we see that $\coprod_{i = 1, \ldots, n} X_i \to Z$ has $\mathcal{P}$.
Now $\{\coprod_{i = 1, \ldots, n} X_i \to X\}$ is an \'etale
covering, hence as $\mathcal{P}$ is \'etale local on the source
we conclude that $X \to Z$ has $\mathcal{P}$ as desired.

\medskip\noindent
Claim 2: (AAA) $\Rightarrow$ (A).
Namely, let $f : X \to Y$, $g : Y \to Z$ be as above with $Z$ affine.
Choose an affine open covering $X = \bigcup X_i$.
As $\mathcal{P}$ is \'etale local on the source we see that
each $f|_{X_i} : X_i \to Y$ has $\mathcal{P}$.
By (AA), which follows from (AAA) according to Claim 1, we see that
$X_i \to Z$ has $\mathcal{P}$ for each $i$.
Since $\{X_i \to X\}$ is an \'etale covering and $\mathcal{P}$ is \'etale
local on the source we conclude that
$X \to Z$ has $\mathcal{P}$.

\medskip\noindent
Claim 3: (AAA) $\Rightarrow$ (-).
Namely, let $f : X \to Y$, $g : Y \to Z$ be as above.
Choose an affine open covering $Z = \bigcup Z_i$.
Set $Y_i = g^{-1}(Z_i)$ and $X_i = f^{-1}(Y_i)$. By
Lemma \href{https://stacks.math.columbia.edu/tag/04QU}{lemma pullback property local target (Tag 04QU)}
each of the morphisms $X_i \to Y_i$ has $\mathcal{P}$.
By (A), which follows from (AAA) according to Claim 2, we see that
$X_i \to Z_i$ has $\mathcal{P}$ for each $i$.
Since $\mathcal{P}$ is local on the target and $X_i = (g \circ f)^{-1}(Z_i)$
we conclude that $X \to Z$ has $\mathcal{P}$.

\medskip\noindent
Thus to prove the lemma it suffices to prove (AAA).
Let $f : X \to Y$ and $g : Y \to Z$ be as above $X, Y, Z$ affine.
Note that an \'etale morphism of affines has universally bounded fibres, see
Morphisms,
Lemma \href{https://stacks.math.columbia.edu/tag/03WS}{lemma etale locally quasi finite (Tag 03WS)} and
Lemma \href{https://stacks.math.columbia.edu/tag/03JA}{lemma locally quasi finite qc source universally bounded (Tag 03JA)}.
Hence we can do induction on the integer $n$ bounding the degree of the fibres
of $Y \to Z$. See
Morphisms, Lemma \href{https://stacks.math.columbia.edu/tag/03WU}{lemma etale universally bounded (Tag 03WU)}
for a description of this integer in the case of an \'etale morphism.
If $n = 1$, then $Y \to Z$ is an open immersion, see
Lemma \ref{lemma-universally-injective-etale-open-immersion},
and the result follows from assumption (3) of the lemma. Assume $n > 1$.

\medskip\noindent
Consider the following commutative diagram
$$
\xymatrix{
X \times_Z Y \ar[d] \ar[r]_{f_Y} &
Y \times_Z Y \ar[d] \ar[r]_-{\text{pr}} &
Y \ar[d] \\
X \ar[r]^f &
Y \ar[r]^g &
Z
}
$$
Note that we have a decomposition into open and closed
subschemes $Y \times_Z Y = \Delta_{Y/Z}(Y) \amalg Y'$, see
Morphisms, Lemma \href{https://stacks.math.columbia.edu/tag/02GE}{lemma diagonal unramified morphism (Tag 02GE)}.
As a base change the degrees of the fibres of the second projection
$\text{pr} : Y \times_Z Y \to Y$ are bounded by $n$, see
Morphisms, Lemma \href{https://stacks.math.columbia.edu/tag/03J7}{lemma base change universally bounded (Tag 03J7)}.
On the other hand, $\text{pr}|_{\Delta(Y)} : \Delta(Y) \to Y$ is
an isomorphism and every fibre has exactly one point.
Thus, on applying
Morphisms, Lemma \href{https://stacks.math.columbia.edu/tag/03WU}{lemma etale universally bounded (Tag 03WU)}
we conclude the degrees of the fibres of the restriction
$\text{pr}|_{Y'} : Y' \to Y$ are bounded by $n - 1$.
Set $X' = f_Y^{-1}(Y')$. Picture
$$
\xymatrix{
X \amalg X' \ar@{=}[d] \ar[r]_-{f \amalg f'} &
\Delta(Y) \amalg Y' \ar@{=}[d] \ar[r] &
Y \ar@{=}[d] \\
X \times_Z Y \ar[r]^{f_Y} &
Y \times_Z Y \ar[r]^-{\text{pr}} &
Y
}
$$
As $\mathcal{P}$ is \'etale local on the target and hence stable under
\'etale base change (see
Lemma \href{https://stacks.math.columbia.edu/tag/04QU}{lemma pullback property local target (Tag 04QU)})
we see that $f_Y$ has $\mathcal{P}$.
Hence, as $\mathcal{P}$ is \'etale local on the source,
$f' = f_Y|_{X'}$ has $\mathcal{P}$. By induction hypothesis
we see that $X' \to Y$ has $\mathcal{P}$.
As $\mathcal{P}$ is local on the source, and
$\{X \to X \times_Z Y, X' \to X \times_Y Z\}$ is an \'etale covering,
we conclude that $\text{pr} \circ f_Y$ has $\mathcal{P}$.
Note that $g \circ f$ can be viewed as a morphism
$g \circ f : X \to g(Y)$. As $\text{pr} \circ f_Y$ is the pullback of
$g \circ f : X \to g(Y)$ via the \'etale covering $\{Y \to g(Y)\}$,
and as $\mathcal{P}$ is \'etale local on the target, we conclude that
$g \circ f : X \to g(Y)$ has property $\mathcal{P}$. Finally, applying
assumption (3) of the lemma once more we conclude that
$g \circ f : X \to Z$ has property $\mathcal{P}$.
\end{proof}
\end{document}
